%
%

\chapter{The naturals, to a thousand places}

Seventy-four constants, every one verified to all thousand places it is published at, and a table
that refuses to be written if that is not true of every row.

\section{Two routes or it does not ship}

A constant computed once is a constant nobody has checked. A single series with a transcribed error
returns a plausible number, and no amount of rereading the code finds it, because the code says what
it was written to say. So every row carries two computations and the gap between them.

\begin{center}
\begin{tabular}{@{}lll@{}}
\toprule
group & first route & second route \\
\midrule
\(\pi\) and its multiples & Chudnovsky, binary splitting & Machin, \(16\arctan\frac{1}{5}-4\arctan\frac{1}{239}\) \\
square roots & integer Newton on a scaled integer & squaring the answer back \\
cube roots & integer Newton, \(n\)th root & raising the answer back \\
logarithms & one \(\operatorname{artanh}\) decomposition & a second sharing no term \\
\(\zeta\) of even integers & \(\pi^k/c\) from Chudnovsky & the same from Machin \\
\(\zeta(3)\) & a BBP-type series, binary scaled & \(\frac{5}{2}\sum(-1)^{k+1}/(k^3\binom{2k}{k})\) \\
\bottomrule
\end{tabular}
\end{center}

\textbf{A root does not get a second series, and pretending otherwise was this tool's first defect.}
It claimed one and then called the same routine twice, which grades nothing at all. No independent
series is needed: \(r^2 \le N < (r+1)^2\) is an exact test with no tolerance in it, and raising an
\(n\)th root back is the same statement. Those rows are stronger than the ones with two series, and not
weaker.

The logarithm identities rest on \(\operatorname{artanh}(1/m) = \frac{1}{2}\ln\frac{m+1}{m-1}\) and
are stated so they can be checked without running anything:

\begin{center}
\begin{tabular}{@{}lll@{}}
\toprule
constant & route A & route B \\
\midrule
\(\ln 3\) & \(2\operatorname{artanh}\frac{1}{2}\) & \(4\operatorname{artanh}\frac{1}{5} + 2\operatorname{artanh}\frac{1}{7}\) \\
\(\ln 5\) & \(2\operatorname{artanh}\frac{1}{9} + 2\ln 2\) & \(2\operatorname{artanh}\frac{1}{4} + 2\operatorname{artanh}\frac{1}{2}\) \\
\(\ln 7\) & \(6\operatorname{artanh}\frac{1}{3} - 2\operatorname{artanh}\frac{1}{15}\) & \(2\operatorname{artanh}\frac{1}{6} + 2\operatorname{artanh}\frac{1}{4} + 2\operatorname{artanh}\frac{1}{2}\) \\
\bottomrule
\end{tabular}
\end{center}

Nothing in that table is compared against digits typed into the source. Digits typed from memory are
how a wrong row ships. The pair therefore decides, and the table consults no reference prefix.

\section{No tolerance anywhere, because none is needed}

Measured at 60, 250 and 1000 places: \textbf{every two-route gap in the file is exactly zero.} Every
threshold is therefore a decision, and not a bound.

An earlier version of these checks allowed a few units of slack. That is a bound chosen by judgment
sitting in front of every row, and such a bound admits precisely the one row that drifted. One check
had it worst. The identity \(5\zeta(4) = 2\zeta(2)^2\) is exact, and the first version of that check
inflated a three-unit truncation by \(10^{\text{digits}}\) before comparing it against a threshold in
unscaled units, a correct identity read as a sixty-digit failure. The repair was then to widen the
threshold, \textbf{which was worse than the arithmetic error that caused it}: computing both sides at
the guarded length and cutting gives exactly zero.

\section{The page guard is a narrowing alarm, not an overflow guard}

Nothing here can overflow. The arithmetic runs on integers with no bound, and it runs on them for that reason.
The failure worth guarding is the silent opposite: a row that keeps its printed width and loses its
meaning in the tail. It prints a thousand digits, the leading nine hundred are right, and nothing
about the output looks wrong.

\begin{center}
\begin{tabular}{@{}lll@{}}
\toprule
 & what happens & how it announces itself \\
\midrule
overflow & the value grows too large & raises or wraps, \textbf{loudly} \\
narrowing & the width holds, the tail is fiction & \textbf{nothing at all} \\
\bottomrule
\end{tabular}
\end{center}

And it compounds: the next constant computed from a narrowed one truncates again, and the place the
loss surfaces is nowhere near where it began.

\textbf{It caught a row that had been shipping as a footnote for as long as the table existed.} The
placement angle agreed to 999 of 1000 places, and the tool printed a line saying a few units in the
last place were the guard truncating. That was true and it was the wrong response. The composition
\(\pi(3-\sqrt5)\) truncates twice, once in the subtraction's operand and once in the product, and
computing it at the guarded width removes the loss entirely for the cost of one extra Machin run.
The footnote taught a reader to expect slack, and the slack was real lost digits.

The guard refuses the write, and does not merely warn. It names every short row. Widening the working
length is the repair; narrowing the request is not, because the request is what a caller relies on.

\textbf{It measures agreement, and not correctness.} Two routes agreeing to \(N\) places means
\(N\) places are verified. Two transcriptions of one wrong formula agree perfectly. Independence of
the routes is a reading of the mathematics, and no tool performs it. This is therefore the second of two
gates, and never the only one.

\section{What is in the table}

\begin{center}
\begin{tabular}{@{}llp{6.5cm}@{}}
\toprule
group & count & note \\
\midrule
\(\pi\), its multiples and powers & 6 & including the placement angle \\
even \(\zeta\) & 3 & \(\zeta(2)\), \(\zeta(4)\), \(\zeta(6)\) \\
\(\zeta(3)\) & 1 & no closed form in \(\pi\): two real series \\
square roots of the naturals to 30 & 24 & perfect squares excluded \\
cube roots of the naturals to 30 & 27 & perfect cubes excluded \\
logarithms & 6 & \(\ln 2,3,5,7,10\) and \(\log_2 10\) \\
the rest & 7 & \(e\), \(\varphi\), the harmonic unit, roots of 2, 3, 5 \\
\bottomrule
\end{tabular}
\end{center}

The table leaves out perfect powers deliberately. \(\sqrt{25}\) is 5, a rational, and a table of irrationals
carrying it would be claiming a thousand places of something with one digit.

\section{What arbitrary precision costs, since the cost is the evidence}

\textbf{A constant that arrives instantly was looked up.} The only proof available from outside that
a number was computed is that it was paid for: a lookup is flat in the digit count and a computation
is not. So the timing mode reports seconds and grades the growth against what each algorithm
predicts from its own shape, never against a fit.

\begin{center}
\begin{tabular}{@{}lrrrrrr@{}}
\toprule
route & 250 & 1000 & 4000 & 16000 & measured & predicted \\
\midrule
\(\pi\), Chudnovsky & 0.0002 & 0.0010 & 0.0091 & 0.0752 & 1.452 & 1.600 \\
\(\pi\), Machin & 0.0002 & 0.0019 & 0.0289 & 0.4207 & 1.830 & 2.000 \\
root of two & 0.0001 & 0.0002 & 0.0015 & 0.0121 & 1.297 & 1.600 \\
\(e\), factorial series & 0.0001 & 0.0005 & 0.0056 & 0.0686 & 1.625 & 2.000 \\
\(\ln 2\), \(\operatorname{artanh}\) & 0.0002 & 0.0020 & 0.0275 & 0.4426 & 1.842 & 2.000 \\
\bottomrule
\end{tabular}
\end{center}

Every row pays for its digits and no row reads near zero in the measured column. The two routes to
\(\pi\) pay differently, and the table carries both for that reason: a series with one division at the end
grows like its multiply, and a series with a division per term grows like the square.

The honest caveat on those exponents is that the interpreter dominates the small end. That inflates
the short timings and flattens the slope, and the true exponents are therefore nearer the predictions than 1.297
suggests.

\textbf{And the spread between routes is the practical shape of the problem.} Chudnovsky gains 14.18
digits a term; the central binomial series for \(\zeta(3)\) gains 0.78. a constant with no rapidly
convergent series costs about eighteen times the terms for the same output, each term a full-precision
operation. The constants with closed forms in \(\pi\) are nearly free once \(\pi\) is paid for, and
the ones without are where the time goes.

\section{Where carrying our own constants buys something, narrowly}

The rule is that this work uses its own constants, and the reason is narrower than it first looks. It
therefore deserves a precise statement, and not an assertion.

Above double precision it is not a preference but a requirement: a standard-library \(\pi\) is a
double, and any computation carrying more than about sixteen digits and reaching for one has silently
capped itself, with the run looking fine throughout.

At double precision a single named constant gains nothing. A correctly rounded library constant and
our own thousand-digit value rounded to a double are the identical bits, and claiming an improvement
there would be false. Measured over fourteen library forms: \textbf{thirteen give bit-identical
doubles.} The one that differs is the composed placement angle, \(\pi(3-\sqrt5)\), which rounds three
times and lands one unit in the last place away. That is the constant every source in the placement
is positioned by.

So the honest case is: required above sixteen digits, free and bit-identical for thirteen of fourteen
at sixteen, and worth one unit in the last place on the one composed constant that matters.
Composition costs, and most compositions in this table cost nothing. That is a measurement,
and its outcome could differ.

\section{Euler--Mascheroni has two real routes}

The first route is Brent--McMillan, a real implementation, and not a series summed casually. The second
independent route is Sweeney's exponential-integral method, a second real implementation. The two fail
differently: Brent--McMillan has all positive terms, and Sweeney loses about 1779 digits to cancellation
before it delivers one. They agree to the full width, and the row carries that agreement.

A row with one unverified number would be worse than no row, because the claim of the table is that
no row is unverified. \textbf{One exception makes the other rows unreadable as a group.}
